La lectura dada determina
\[
v=\frac{0.8}{\sqrt3}\approx0.461880\quad\mathrm{m/s}.
\]
En general,
\[
x(\tau)=0.30\cos(2\tau)+\frac v2\sin(2\tau).
\]
Si \(\sin(2\tau)\ne0\), hay una única velocidad,
\[
v=\frac{2[z-0.30\cos(2\tau)]}{\sin(2\tau)},\qquad
\Delta v=\frac{2\eta}{\sin(2\tau)}.
\]
En \(\tau=n\pi/2\), \(n=1,2,\ldots\), la lectura \(z=0.30(-1)^n\) es compatible con cualquier \(v\); cualquier otra lectura es incompatible. Cerca de esos instantes, el factor que amplifica los cambios de lectura crece sin cota.
